% ============================================================
%  Glosario
%  TT 2026-B066 | ESCOM-IPN
% ============================================================

\chapter*{Glosario}
\addcontentsline{toc}{chapter}{Glosario}
\markboth{Glosario}{Glosario}

\setlength{\LTpre}{6pt}
\setlength{\LTpost}{0pt}

\begin{longtable}{@{} p{4.3cm} p{11.3cm} @{}}

  \toprule
  \textbf{Término} & \textbf{Definición} \\
  \midrule
  \endfirsthead

  \multicolumn{2}{@{}l}{\small\itshape (continúa de la página anterior)} \\[4pt]
  \toprule
  \textbf{Término} & \textbf{Definición} \\
  \midrule
  \endhead

  \midrule
  \multicolumn{2}{r@{}}{\small\itshape (continúa en la página siguiente)} \\
  \endfoot

  \bottomrule
  \endlastfoot

  % ── A ──────────────────────────────────────────────────────────────
  \textbf{Antidepresivo}
  & Fármaco que reduce los síntomas depresivos modulando la disponibilidad de
    neurotransmisores monoaminérgicos (serotonina, noradrenalina o dopamina) en la
    hendidura sináptica. En el contexto del FST, su efecto se mide como reducción
    del tiempo de inmovilidad. \\[5pt]

  \textbf{API REST}
  & Interfaz de programación de aplicaciones que sigue las restricciones
    arquitectónicas REST (\textit{Representational State Transfer}). En el
    sistema propuesto, la API Flask expone los recursos del backend (experimentos,
    videos, resultados) mediante peticiones HTTP estándar (GET, POST, PATCH, DELETE)
    e intercambia datos en formato JSON. \\[5pt]

  \textbf{Aprendizaje supervisado}
  & Paradigma de aprendizaje automático en el que el modelo se entrena con
    pares entrada-etiqueta proporcionados por un experto. En el proyecto, el
    clasificador de conductas se entrena con fotogramas etiquetados manualmente
    por el colaborador clínico. \\[5pt]

  % ── B ──────────────────────────────────────────────────────────────
  \textbf{Behavioral despair}
  & (Desesperanza conductual.) Estado conductual inducido por exposición a un
    estresor inescapable —como el agua en el FST— que se manifiesta como
    inmovilidad progresiva. Postulado por Porsolt (1977) como modelo
    traslacional de depresión. \\[5pt]

  \textbf{Bounding box}
  & (Caja delimitadora.) Rectángulo mínimo que encierra a un objeto detectado
    en una imagen. El clasificador ResNet recibe el recorte de imagen definido
    por el bounding box de la rata, normalizado a $224\times224$ píxeles. \\[5pt]

  \textbf{ByteTrack}
  & Algoritmo de seguimiento multi-objeto que asocia detecciones entre fotogramas
    usando IoU y un filtro de Kalman. Se utiliza junto con YOLOv8 para mantener
    la identidad de cada rata a lo largo del video. \\[5pt]

  % ── C ──────────────────────────────────────────────────────────────
  \textbf{CLAHE}
  & (\textit{Contrast Limited Adaptive Histogram Equalization}.) Técnica de
    mejora de contraste local que divide la imagen en teselas y aplica
    ecualización de histograma con un límite de amplificación. Se aplica como
    preprocesamiento cuando el contraste global del video es insuficiente. \\[5pt]

  \textbf{CNN}
  & (\textit{Convolutional Neural Network} / Red Neuronal Convolucional.)
    Arquitectura de red profunda que aplica filtros convolucionales para extraer
    características espaciales jerárquicas de las imágenes. El clasificador de
    conductas del sistema es una CNN (ResNet-18) reentrenada con datos del FST. \\[5pt]

  \textbf{Cola de tareas}
  & Estructura de datos FIFO que desacopla la recepción de solicitudes de análisis
    del procesamiento efectivo. El worker consulta la cola periódicamente y toma
    un trabajo a la vez, lo que permite procesar múltiples videos de forma
    secuencial sin bloquear el servidor web. \\[5pt]

  \textbf{Conducta de escalamiento}
  & (\textit{Climbing / escape attempts}.) Comportamiento en el FST caracterizado
    por movimientos verticales enérgicos de las patas anteriores contra las
    paredes del cilindro. Se asocia con la activación del sistema noradrenérgico. \\[5pt]

  \textbf{Conducta de nado activo}
  & (\textit{Swimming}.) Movimientos horizontales del cuerpo y las extremidades
    que desplazan al animal por el cilindro. Se asocia con la actividad del
    sistema serotoninérgico y se incrementa con antidepresivos ISRS. \\[5pt]

  \textbf{Contorno}
  & Curva que delimita los bordes de un objeto en una imagen binarizada. En el
    pipeline, los contornos obtenidos tras la umbralización de Otsu sirven para
    detectar la silueta de la rata y calcular su bounding box. \\[5pt]

  \textbf{Corticosterona}
  & Principal glucocorticoide en roedores, secretado por las glándulas suprarrenales
    en respuesta al estrés. Sus niveles plasmáticos se elevan significativamente
    durante el FST y constituyen un biomarcador del eje hipotálamo-hipófisis-adrenal. \\[5pt]

  \textbf{CSRT}
  & (\textit{Channel and Spatial Reliability Tracker}.) Algoritmo de seguimiento
    por filtro de correlación discriminativo que pondera las regiones de la imagen
    con mayor fiabilidad espacial. Se utiliza como respaldo cuando la segmentación
    por sustracción de fondo falla en un fotograma. \\[5pt]

  % ── D ──────────────────────────────────────────────────────────────
  \textbf{Data augmentation}
  & (Aumento de datos.) Técnica que incrementa artificialmente el tamaño del
    conjunto de entrenamiento aplicando transformaciones geométricas (rotación,
    volteo, zoom) y fotométricas (brillo, contraste) a las imágenes originales,
    reduciendo el riesgo de sobreajuste. \\[5pt]

  \textbf{Detección de objetos}
  & Tarea de visión por computadora que consiste en localizar y clasificar
    simultáneamente múltiples instancias de objetos en una imagen. En el proyecto
    se emplea YOLOv8 para detectar las ratas dentro de sus respectivas ROIs. \\[5pt]

  \textbf{Docker}
  & Plataforma de contenedorización que empaqueta una aplicación junto con todas
    sus dependencias en unidades aisladas y reproducibles. El sistema se
    despliega como cuatro contenedores orquestados: frontend, backend, worker y
    base de datos. \\[5pt]

  \textbf{Dopaminérgico}
  & Relativo al sistema de neurotransmisión mediado por dopamina. Fármacos que
    potencian la actividad dopaminérgica pueden modificar el patrón de conductas
    en el FST, aunque su efecto es menos selectivo que el de los ISRS sobre el
    tiempo de nado activo. \\[5pt]

  % ── E ──────────────────────────────────────────────────────────────
  \textbf{EMA}
  & (\textit{Exponential Moving Average} / Promedio móvil exponencial.)
    Filtro de suavizado que asigna mayor peso a las observaciones recientes.
    En el pipeline se utiliza para estabilizar la estimación de la línea de
    agua y reducir el efecto de reflejos transitorios en la detección. \\[5pt]

  \textbf{Entropía cruzada}
  & (\textit{Cross-entropy loss}.) Función de pérdida estándar para clasificación
    multiclase que mide la discrepancia entre la distribución de probabilidad
    predicha por el modelo y la distribución real (etiquetas one-hot). Se
    minimiza durante el entrenamiento del clasificador de conductas. \\[5pt]

  % ── F ──────────────────────────────────────────────────────────────
  \textbf{F1-score}
  & Media armónica entre precisión y recall: $F_1 = 2 \cdot \frac{P \cdot R}{P + R}$.
    Métrica balanceada que penaliza tanto los falsos positivos como los falsos
    negativos; se reporta por clase (inmovilidad, nado activo, escalamiento)
    para evaluar el clasificador. \\[5pt]

  \textbf{Frame}
  & (Cuadro de video.) Imagen estática individual dentro de una secuencia de
    video. El análisis del FST opera cuadro a cuadro (\textit{frame-by-frame}):
    cada fotograma es clasificado independientemente y los tiempos de conducta
    se acumulan en segundos sumando los cuadros clasificados en cada categoría. \\[5pt]

  \textbf{FST}
  & (\textit{Forced Swim Test} / Prueba de Nado Forzado.) Paradigma experimental
    propuesto por Porsolt et al.\ (1977) en el que una rata se coloca en un
    cilindro con agua de la que no puede escapar. La duración de la inmovilidad
    en la segunda sesión (día 2) es el indicador primario de desesperanza
    conductual y sensibilidad a antidepresivos. \\[5pt]

  % ── G ──────────────────────────────────────────────────────────────
  \textbf{Gold standard}
  & (Estándar de referencia.) Conjunto de etiquetas considerado como la verdad
    absoluta para la evaluación del sistema. En el proyecto, lo constituyen las
    anotaciones manuales cuadro a cuadro realizadas por el experto del laboratorio
    colaborador, contra las cuales se comparan las predicciones del modelo. \\[5pt]

  % ── H ──────────────────────────────────────────────────────────────
  \textbf{Hiperparámetro}
  & Parámetro cuyo valor se fija antes del entrenamiento y no se aprende de los
    datos (e.g., tasa de aprendizaje, tamaño del lote, número de épocas,
    umbral de confianza del clasificador). Se ajusta mediante búsqueda en
    cuadrícula o validación cruzada. \\[5pt]

  % ── I ──────────────────────────────────────────────────────────────
  \textbf{ImageNet}
  & Base de datos de más de 14 millones de imágenes anotadas en más de 21{,}000
    categorías. Los pesos del modelo ResNet-18 preentrenados en ImageNet se usan
    como punto de partida para el transfer learning, aprovechando características
    visuales generales antes de la especialización en fotogramas del FST. \\[5pt]

  \textbf{Inmovilidad}
  & (\textit{Immobility}.) Conducta en el FST caracterizada por la ausencia de
    movimientos propulsivos; el animal flota con movimientos mínimos para
    mantener la cabeza sobre el agua. Es la medida primaria de la desesperanza
    conductual y el indicador de referencia para cuantificar el efecto
    antidepresivo. \\[5pt]

  \textbf{IoU}
  & (\textit{Intersection over Union}.) Métrica de solapamiento entre dos
    rectángulos (generalmente una detección y su verdad de referencia):
    $\text{IoU} = \frac{|A \cap B|}{|A \cup B|}$. En el sistema se usa para
    validar la consistencia entre detecciones consecutivas durante el seguimiento
    de las ratas. \\[5pt]

  \textbf{ISRS}
  & (Inhibidor Selectivo de la Recaptación de Serotonina.) Clase farmacológica
    que bloquea el transportador SERT, incrementando la disponibilidad sináptica
    de serotonina. En el FST aumentan específicamente el tiempo de nado activo
    sin afectar de forma significativa el escalamiento. Fluoxetina y sertralina
    son ejemplos representativos. \\[5pt]

  % ── J ──────────────────────────────────────────────────────────────
  \textbf{JWT}
  & (\textit{JSON Web Token}.) Estándar abierto (RFC 7519) para la transmisión
    segura de información entre partes como un objeto JSON firmado digitalmente.
    El backend emite un JWT al iniciar sesión; el frontend lo adjunta en el
    encabezado \texttt{Authorization} de cada petición subsecuente para
    autenticar al usuario. \\[5pt]

  % ── K ──────────────────────────────────────────────────────────────
  \textbf{Kappa de Cohen}
  & Coeficiente estadístico ($\kappa$) que mide el grado de concordancia entre
    dos evaluadores descontando el acuerdo esperado por azar. Valores $\kappa > 0.80$
    se consideran acuerdo casi perfecto. Se reporta para cuantificar la concordancia
    entre el sistema automático y el experto en la clasificación cuadro a cuadro. \\[5pt]

  \textbf{KCF}
  & (\textit{Kernelized Correlation Filter}.) Algoritmo de seguimiento basado en
    correlación en el dominio de Fourier con kernel RBF, evaluado como alternativa
    al CSRT durante el diseño del módulo de rastreo. El CSRT fue seleccionado
    por su mayor precisión en presencia de oclusión parcial. \\[5pt]

  % ── M ──────────────────────────────────────────────────────────────
  \textbf{MAE}
  & (\textit{Mean Absolute Error} / Error Absoluto Medio.) Métrica de error de
    regresión definida como $\text{MAE} = \frac{1}{n}\sum_{i=1}^{n}|y_i - \hat{y}_i|$.
    En la evaluación del sistema se calcula entre los tiempos de conducta
    estimados automáticamente y los tiempos del gold standard, expresado
    en segundos. \\[5pt]

  \textbf{Modelo de fondo}
  & (\textit{Background model}.) Representación estadística de la escena sin el
    objeto de interés, construida en el pipeline mediante la mediana de los primeros
    cuadros del video. Permite detectar la rata por diferencia absoluta entre
    cada fotograma y el fondo. \\[5pt]

  % ── N ──────────────────────────────────────────────────────────────
  \textbf{Noradrenérgico}
  & Relativo al sistema de neurotransmisión mediado por noradrenalina (norepinefrina).
    Los antidepresivos noradrenérgicos (e.g., desipramina) aumentan
    preferentemente el tiempo de escalamiento en el FST, a diferencia de los
    serotoninérgicos que aumentan el nado activo. \\[5pt]

  % ── O ──────────────────────────────────────────────────────────────
  \textbf{ORM}
  & (\textit{Object-Relational Mapping}.) Técnica que abstrae las operaciones de
    base de datos relacional mediante objetos del lenguaje de programación. El
    backend utiliza SQLAlchemy como ORM para interactuar con PostgreSQL sin
    escribir SQL directamente en la lógica de negocio. \\[5pt]

  \textbf{Otsu, umbralización de}
  & Método automático de binarización de imágenes que calcula el umbral óptimo
    minimizando la varianza intraclase entre los píxeles de fondo y primer plano.
    En el pipeline se aplica al mapa de diferencias para segmentar la silueta de
    la rata del fondo. \\[5pt]

  \textbf{Overfitting}
  & (Sobreajuste.) Fenómeno en el que un modelo aprende el ruido del conjunto de
    entrenamiento en lugar de los patrones generalizables, resultando en alta
    precisión en entrenamiento pero baja en datos nuevos. Se mitiga con data
    augmentation, dropout y validación cruzada. \\[5pt]

  % ── P ──────────────────────────────────────────────────────────────
  \textbf{Pipeline de análisis}
  & Secuencia de etapas de procesamiento encadenadas en las que la salida de
    cada etapa es la entrada de la siguiente. En el sistema, el pipeline comprende:
    preprocesamiento, detección de ROIs, seguimiento con YOLO, y clasificación
    con ResNet, ejecutado por el worker de forma asíncrona. \\[5pt]

  \textbf{PostgreSQL}
  & Sistema de gestión de bases de datos relacional de código abierto con soporte
    para tipos de datos avanzados y transacciones ACID. Es la base de datos
    primaria del sistema B066, que almacena usuarios, experimentos, videos,
    resultados de conducta y notificaciones. \\[5pt]

  \textbf{Precisión}
  & (\textit{Precision}.) Fracción de las predicciones positivas que son
    verdaderos positivos: $P = \frac{VP}{VP + FP}$. Una precisión alta indica
    pocos falsos positivos en la clasificación de cada conducta. \\[5pt]

  \textbf{Prueba de nado forzado}
  & Véase \textbf{FST}. Protocolo de dos sesiones (día 1: 15 min de nado de
    habituación; día 2: 5 min de prueba) durante las cuales se registran en
    video los tiempos de inmovilidad, nado activo y escalamiento de cada animal. \\[5pt]

  % ── R ──────────────────────────────────────────────────────────────
  \textbf{Recall}
  & (Sensibilidad / \textit{Sensitivity}.) Fracción de los casos positivos reales
    que son correctamente identificados: $R = \frac{VP}{VP + FN}$. Un recall alto
    indica que el sistema detecta la mayoría de los fotogramas de cada conducta
    presentes en el video. \\[5pt]

  \textbf{ResNet}
  & (\textit{Residual Network}.) Familia de redes neuronales convolucionales con
    conexiones residuales (atajos) que permiten entrenar arquitecturas muy
    profundas sin degradación del gradiente. Se usa ResNet-18 como clasificador
    de conductas por su balance entre capacidad expresiva y tiempo de inferencia. \\[5pt]

  \textbf{ROI}
  & (\textit{Region of Interest} / Región de interés.) Subregión rectangular del
    fotograma que delimita el área de análisis de un animal específico. En el
    sistema, cada cilindro del FST corresponde a una ROI detectada automáticamente
    al inicio del pipeline. \\[5pt]

  % ── S ──────────────────────────────────────────────────────────────
  \textbf{Segmentación}
  & Proceso que asigna una etiqueta (clase) a cada píxel de una imagen. En el
    contexto del pipeline, la segmentación por sustracción de fondo separa los
    píxeles correspondientes a la rata de los píxeles del fondo estático para
    obtener su contorno y bounding box. \\[5pt]

  \textbf{Serotoninérgico}
  & Relativo al sistema de neurotransmisión mediado por serotonina (5-HT). Los
    antidepresivos serotoninérgicos (ISRS) incrementan específicamente el tiempo
    de nado activo en el FST sin modificar significativamente el escalamiento ni
    la actividad locomotora general. \\[5pt]

  \textbf{Softmax}
  & Función de activación que transforma un vector de puntuaciones en una
    distribución de probabilidad: $\sigma(z)_i = e^{z_i}/\sum_j e^{z_j}$.
    La capa final del clasificador ResNet aplica softmax para obtener la
    probabilidad de cada conducta (inmovilidad, nado activo, escalamiento). \\[5pt]

  \textbf{Sustracción de fondo}
  & (\textit{Background subtraction}.) Técnica de segmentación que identifica
    los píxeles que cambian respecto al modelo de fondo estático. Es el método
    primario del pipeline para detectar la presencia y posición de la rata
    en cada cuadro del video. \\[5pt]

  % ── T ──────────────────────────────────────────────────────────────
  \textbf{Tracker}
  & (Rastreador.) Algoritmo que mantiene la identidad de un objeto entre
    fotogramas consecutivos de un video sin necesidad de re-detectarlo en cada
    cuadro. En el sistema se utiliza CSRT como respaldo cuando la segmentación
    por sustracción de fondo falla. \\[5pt]

  \textbf{Transfer learning}
  & (Aprendizaje por transferencia.) Técnica que reutiliza los pesos de un
    modelo entrenado en una tarea o conjunto de datos extenso (ImageNet) como
    punto de partida para una tarea relacionada con datos más limitados. Permite
    entrenar el clasificador de conductas con menos fotogramas etiquetados. \\[5pt]

  % ── W ──────────────────────────────────────────────────────────────
  \textbf{Worker}
  & Proceso independiente del servidor web que consume trabajos de la cola de
    tareas y ejecuta el pipeline de análisis completo. Al correr en un contenedor
    Docker separado, permite que el backend siga respondiendo peticiones HTTP
    mientras se procesa un video de larga duración. \\[5pt]

  % ── Y ──────────────────────────────────────────────────────────────
  \textbf{YOLO}
  & (\textit{You Only Look Once}.) Familia de detectores de objetos en tiempo
    real basados en una sola pasada por la red para predecir bounding boxes y
    clases simultáneamente. En el sistema se emplea YOLOv8 para detectar y
    localizar las ratas dentro de las ROIs de cada cilindro. \\[5pt]

\end{longtable}
